\chapter{Law Without Referencing the Law}
\label{short:law}

Sometimes naming the rule that protects a person may expose that protection
to alteration or retaliation. The original question imagines public rules
and a hidden event: could a decision establish compliance, or a violation,
without exposing the event or identifying the decisive rule?

Secrecy may preserve an advantage that freedom needs. It may also leave a
person unable to understand or contest a judgment. Who fixes the rules? Who
can inspect the decision, demand reasons, or appeal?

The chapter offers this tension for investigation. No judicial protocol has
been established. Any construction would have to confront the power it gives
to those who judge.

\shortlimit{A hidden verification process cannot establish due process by itself.}
\companionref{13}{chapter.13}{A design intuition; authority, privacy, reasons, and appeal remain open.}
